\section*{Chapitre \ref{ch:b3:advanced-nmr} --- La RMN approfondie~: impulsions, relaxation et 2D}

\begin{solution}{exo:b3:advanced-nmr:1}
$\nu_0 = (\gamma/2\pi)B_0$~: \ce{^1H} \qty{600.4}{MHz}, \ce{^{13}C} \qty{151.0}{MHz}, \ce{^{19}F}
\qty{565.1}{MHz}.
\end{solution}

\begin{solution}{exo:b3:advanced-nmr:2}
$h\nu_0/2kT = 6.626\times10^{-34} \times 400.2\times10^6/(2 \times 1.381\times10^{-23} \times 300)
= \num{3.2e-5}$. À \qty{4}{K}, il est 75 fois plus grand, \num{2.4e-3}
(l'approximation reste valable).
\end{solution}

\begin{solution}{exo:b3:advanced-nmr:3}
$\theta = \gamma B_1t_p = \pi/2$ en \qty{10}{\micro s}~: $\gamma B_1/2\pi = 1/(4 \times
\qty{10}{\micro s}) = \qty{25}{kHz}$. Une impulsion de $180^\circ$ dure
\qty{20}{\micro s}~; \qty{5}{\micro s}
donnent $45^\circ$.
\end{solution}

\begin{solution}{exo:b3:advanced-nmr:4}
$T_2 = 1/(\pi \times 3.2) = \qty{0.10}{s}$.
\end{solution}

\begin{solution}{exo:b3:advanced-nmr:5}
$(10.71/42.58)^3 = 0.0159$, multiplié par 0.011~: \num{1.75e-4}, environ
1/5700 du signal du proton. Comme S/B $\propto\sqrt N$, un même rapport
S/B exige $5700^2 \approx \num{3e7}$ fois plus d'acquisitions ---
impossible, et c'est pourquoi les spectres de ${}^{13}$C utilisent des
échantillons plus concentrés, le découplage et le transfert de
polarisation.
\end{solution}

\begin{solution}{exo:b3:advanced-nmr:6}
$T_1 = 1.39/\ln2 = \qty{2.0}{s}$~; délai de recyclage $\ge 5T_1 = \qty{10}{s}$.
\end{solution}

\begin{solution}{exo:b3:advanced-nmr:7}
\omterm{def:b3:advanced-nmr:experiments}{COSY}~: un seul \omterm{def:b3:advanced-nmr:two-d}{pic de corrélation}, entre le quadruplet \ce{CH2} et le
triplet \ce{CH3} du groupe éthyle~; le singulet \ce{CH3CO} n'est couplé à
rien. \omterm{def:b3:advanced-nmr:experiments}{HMBC} depuis le méthyle singulet~: le carbone du carbonyle (deux
liaisons) et le carbone \ce{CH2} (trois liaisons).
\end{solution}

\begin{solution}{exo:b3:advanced-nmr:8}
$k_c = \pi \times 50/\sqrt2 = \qty{111}{s^{-1}}$. $\Delta G^\ddagger = 8.314 \times 330 \times
\ln(6.88\times10^{12}/111) = \qty{68}{kJ/mol}$.
\end{solution}

\begin{solution}{exo:b3:advanced-nmr:9}
NOE $\propto r^{-6}$~: $r_b = 0.30 \times 4^{-1/6} = \qty{0.24}{nm}$.
\end{solution}

\begin{solution}{exo:b3:advanced-nmr:10}
Dans le \omterm{def:b3:advanced-nmr:magnetisation}{référentiel tournant}, à la résonance, il ne reste que
$\mathbf B_1 = B_1\mathbf e_{x'}$~:
$\dot M_{x'} = 0$, $\dot M_{y'} = \gamma B_1M_z$, $\dot M_z = -\gamma B_1M_{y'}$ (composantes de
$\gamma\mathbf M\times\mathbf B_1$). Donc $M_{x'}$ est constante et
$(M_{y'}, M_z)$ tourne à la vitesse angulaire
$\gamma B_1$~: au bout de $t_p$, de l'angle $\gamma B_1t_p$ autour de $x'$.
\end{solution}

\begin{solution}{exo:b3:advanced-nmr:11}
Un spin soumis à un champ un peu plus intense gagne de la phase à une
vitesse supplémentaire constante pendant
$\tau$~; l'impulsion de $180^\circ$ inverse sa phase, et il en perd
exactement autant pendant le second intervalle $\tau$~: tous ces spins se
réalignent à $2\tau$. Les interactions spin--spin aléatoires fluctuent
dans le temps et ne se répètent pas dans le second intervalle~: le
déphasage qu'elles causent n'est pas défait. L'amplitude de l'écho
décroît avec le vrai $T_2$.
\end{solution}

\begin{solution}{exo:b3:advanced-nmr:12}
Éthanoate d'éthyle, \ce{CH3COOCH2CH3}~: quadruplet \ce{OCH2} (4.12),
singulet \ce{CH3CO} (2.05),
triplet \ce{CH3}~; les protons du quadruplet voient le carbone du
carbonyle à trois liaisons, à travers l'oxygène de l'ester. Le propanoate
de méthyle, \ce{CH3CH2COOCH3}, présenterait son quadruplet vers 2.3 ppm
(sur un carbone vers 27 ppm en \omterm{def:b3:advanced-nmr:experiments}{HSQC}, et non 60) et un singulet
\ce{OCH3} vers 3.7 ppm dont le pic \omterm{def:b3:advanced-nmr:experiments}{HMBC} vers le carbonyle traverse
l'oxygène.
\end{solution}

\begin{solution}{pb:b3:advanced-nmr:1}
\textbf{1.} $(2 \times 6 + 2 - 12)/2 = 1$~: un C=O.
\textbf{2.} 173.6 ppm~: le carbonyle de l'ester~; 60.1 ppm~: le \ce{OCH2}.
\textbf{3.} Cinq~; DEPT-135~: \ce{CH2} négatifs (60.1, 36.3, 18.6),
\ce{CH3} positifs (14.3,
13.7), le carbonyle absent.
\textbf{4.} \qty{400.2}{MHz} et \qty{100.7}{MHz}.
\textbf{5.} $3J = \qty{21.6}{Hz} = \qty{0.054}{ppm}$.
\textbf{6.} Les multiplicités montrent quels groupes sont voisins à
l'intérieur de chaque fragment, mais non comment les fragments se
raccordent à travers l'oxygène et le carbonyle.
\textbf{7.} 4.13--1.25~; 2.28--1.66~; 1.66--0.95 (et leurs symétriques).
\textbf{8.} \ce{OCH2CH3} (4.13, 1.25) et \ce{CH2CH2CH3} (2.28, 1.66, 0.95).
\textbf{9.} Les protons \ce{OCH2} et \ce{CH2C=O} sont séparés par
l'oxygène et le carbone du carbonyle~: cinq liaisons, pas de couplage
résolu.
\textbf{10.} 1.66 ppm, entre \ce{CH2} et \ce{CH3}~: couplé à
2 + 3 = 5 protons avec des $J$
voisines, un sextuplet (un multiplet).
\textbf{11.} Des méthyles couplés supposeraient deux \ce{CH3} sur des
carbones voisins (un motif \ce{CH3-CH3})~: impossible ici.
\textbf{12.} \ce{-O-CH2-CH3} et \ce{-CH2-CH2-CH3}, plus le \ce{C=O}.
\textbf{13.} 4.13/60.1~; 2.28/36.3~; 1.66/18.6~; 1.25/14.3~; 0.95/13.7.
\textbf{14.} 14.3 (\ce{CH3} de l'éthyle) et 13.7 ppm (\ce{CH3} du
butanoyle), distants de \qty{0.6}{ppm} seulement.
\textbf{15.} De 4.13 ppm vers 173.6 ppm~: H--C--O--C, trois liaisons.
\textbf{16.} 2.28 ppm (deux liaisons) et 1.66 ppm (trois liaisons) vers
173.6 ppm.
\textbf{17.} \ce{CH3CH2CH2C(=O)OCH2CH3}, le butanoate d'éthyle. Le
propanoate de propyle placerait un triplet
\ce{OCH2} (et non un quadruplet) vers 4.0 ppm et un quadruplet vers
2.3 ppm.
\textbf{18.} Les couplages à quatre liaisons sont en général inférieurs à
1 Hz~; le délai \omterm{def:b3:advanced-nmr:experiments}{HMBC}, réglé pour 5 à 10 Hz, ne les sélectionne pas.
\textbf{19.} Les carbones relaxent à des vitesses différentes et reçoivent
du découplage des exaltations Overhauser nucléaires différentes.
\textbf{20.} $1 - \eu^{-5/20} = 0.22$~: le carbonyle est sous-représenté
d'un facteur 4.
\textbf{21.} $1 - \eu^{-5} = 0.993$.
\textbf{22.} Le NOE dû au découplage des protons, différent pour chaque
carbone~; on le supprime par le découplage inverse alterné (découpleur
actif seulement pendant l'acquisition).
\textbf{23.} $256 \times \qty{100}{s} = \qty{25600}{s}$, environ 7 heures.
\textbf{24.} \textbf{$5T_1$ du carbone du carbonyle~: \qty{100}{s} entre
deux acquisitions.}
\end{solution}
